\documentclass[11pt, a4paper]{article}

\usepackage{amsmath, amssymb, amsthm}
\usepackage{geometry}
\usepackage{booktabs}
\usepackage{array}
\usepackage{parskip}
\usepackage{microtype}
\usepackage{hyperref}
\usepackage{enumitem}
\usepackage{xcolor}
\usepackage{mdframed}

\geometry{margin=2.5cm, top=2.8cm, bottom=2.8cm}

\hypersetup{
    colorlinks=true,
    linkcolor=black,
    urlcolor=black,
}

% Exercise environment
\newcounter{exercise}[section]
\renewcommand{\theexercise}{\thesection.\arabic{exercise}}

\newenvironment{exercise}[1][]{%
    \refstepcounter{exercise}%
    \par\medskip\noindent%
    \textbf{Exercise~\theexercise}%
    \ifx\relax#1\relax\else\ ---\ #1\fi\par\smallskip%
}{%
    \par\medskip%
}

% Motivation box
\newmdenv[
    linewidth=0.6pt,
    linecolor=black!40,
    backgroundcolor=black!3,
    innerleftmargin=10pt,
    innerrightmargin=10pt,
    innertopmargin=6pt,
    innerbottommargin=6pt,
    skipabove=\medskipamount,
    skipbelow=\medskipamount,
]{motivbox}

% Solution environment
\newcounter{solution}[section]
\renewcommand{\thesolution}{\thesection.\arabic{solution}}

\newenvironment{solution}[1][]{%
    \refstepcounter{solution}%
    \par\medskip\noindent%
    \textbf{Solution~\thesolution}%
    \ifx\relax#1\relax\else\ ---\ #1\fi\par\smallskip%
}{%
    \par\medskip%
}

\newcommand{\code}[1]{\texttt{#1}}
\renewcommand{\vec}[1]{\mathbf{#1}}

\title{\textbf{Mathematics for Machine Learning --- Chapter 2}\\[0.4em]
\large Pencil-and-Paper Exercises}
\author{}
\date{}

\begin{document}

\maketitle
\thispagestyle{empty}

\begin{abstract}
\noindent
\textbf{Who this is for.} Someone working through MML Chapter 2 (Linear Algebra, \S2.1--2.3
and \S2.7) alongside mechanistic interpretability study. These exercises cover systems of
equations, matrix properties (rank, transpose, inverse), and linear mappings --- the
foundations that make transformer weight matrices readable.

\textbf{How to use this.} Work through each section in order. Do the exercises by hand
before checking the solutions starting on page~\pageref{sec:solutions}.
\end{abstract}

\tableofcontents
\newpage

% ──────────────────────────────────────────────────────────────────────────────
\section{Systems of Linear Equations}
% ──────────────────────────────────────────────────────────────────────────────

\begin{motivbox}
\textbf{Why this matters.}
A system of linear equations $A\vec{x} = \vec{b}$ is the fundamental operation underlying
almost everything in a transformer. The question ``what input $\vec{x}$ produces output
$\vec{b}$ through matrix $A$?'' is exactly what you are asking when you reverse-engineer a
weight matrix. Gaussian elimination --- the row-reduction method for solving systems --- also
builds intuition for rank and linear dependence, which reappear constantly in circuit
analysis.
\end{motivbox}

\begin{exercise}[Solving a $2 \times 2$ system]
Solve the following system using substitution or elimination:
\[
    2x + y = 5, \qquad x - y = 1
\]
\begin{enumerate}[label=(\alph*)]
    \item Write this as an augmented matrix $[A \mid \vec{b}]$.
    \item Apply row operations to reach row echelon form.
    \item Read off the solution.
\end{enumerate}
\end{exercise}

\begin{exercise}[Gaussian elimination on a $3 \times 3$ system]
Solve:
\[
    \begin{bmatrix} 1 & 2 & 1 \\ 2 & 1 & 3 \\ 0 & 1 & 1 \end{bmatrix}
    \begin{bmatrix} x_1 \\ x_2 \\ x_3 \end{bmatrix}
    =
    \begin{bmatrix} 6 \\ 9 \\ 3 \end{bmatrix}
\]
\begin{enumerate}[label=(\alph*)]
    \item Write the augmented matrix $[A \mid \vec{b}]$.
    \item Eliminate $x_1$ from rows 2 and 3 using row operations.
    \item Eliminate $x_2$ from row 3.
    \item Back-substitute to find $x_3$, then $x_2$, then $x_1$.
\end{enumerate}
\end{exercise}

\begin{exercise}[Counting solutions]
For each augmented matrix below (already in row echelon form), state whether the system
has \textbf{no solution}, \textbf{one solution}, or \textbf{infinitely many solutions},
and give a brief reason.
\begin{enumerate}[label=(\alph*)]
    \item $\begin{bmatrix} 1 & 2 & \mid & 3 \\ 0 & 1 & \mid & 1 \end{bmatrix}$
    \item $\begin{bmatrix} 1 & 2 & \mid & 3 \\ 0 & 0 & \mid & 5 \end{bmatrix}$
    \item $\begin{bmatrix} 1 & 2 & 3 & \mid & 6 \\ 0 & 1 & 1 & \mid & 3 \\ 0 & 0 & 0 & \mid & 0 \end{bmatrix}$
\end{enumerate}
\end{exercise}

\begin{exercise}[Homogeneous systems]
A \emph{homogeneous} system has the form $A\vec{x} = \vec{0}$.
\begin{enumerate}[label=(\alph*)]
    \item Why does a homogeneous system always have at least one solution?
    \item Find the general solution to:
    \[
        \begin{bmatrix} 1 & 2 & -1 \\ 2 & 4 & -2 \end{bmatrix}
        \begin{bmatrix} x_1 \\ x_2 \\ x_3 \end{bmatrix}
        = \begin{bmatrix} 0 \\ 0 \end{bmatrix}
    \]
    \item The solution set you found is a \emph{subspace}. What is its dimension?
          What does that tell you about the null space of this matrix?
\end{enumerate}
\end{exercise}

\begin{exercise}[Geometric interpretation of solutions]
Consider a system $A\vec{x} = \vec{b}$ where $A$ is $2 \times 2$ and the two rows of the
augmented matrix represent two lines in the plane.
\begin{enumerate}[label=(\alph*)]
    \item Sketch (roughly) three scenarios: (i) the lines intersect at a point,
          (ii) the lines are parallel but distinct, (iii) the lines are identical.
          Match each to: unique solution, no solution, infinitely many solutions.
    \item Solve the system $x + y = 3$, $2x + 2y = 6$. Which case is this?
          Describe the solution set geometrically.
    \item Solve $x + y = 3$, $2x + 2y = 7$. Which case is this?
\end{enumerate}
\end{exercise}

\begin{exercise}[Back-substitution practice]
The following matrix is already in row echelon form. Back-substitute to find the solution:
\[
    \left[\begin{array}{ccc|c}
    2 & -1 & 3 & 7 \\
    0 & 3 & -1 & 4 \\
    0 & 0 & 2 & 6
    \end{array}\right]
\]
Show each step explicitly.
\end{exercise}

% ──────────────────────────────────────────────────────────────────────────────
\section{Matrix Properties}
% ──────────────────────────────────────────────────────────────────────────────

\begin{motivbox}
\textbf{Why this matters.}
Transformer weight matrices have specific structural properties --- the QK matrix is low-rank
(at most \texttt{d\_head} rank), the attention pattern matrix is square, and composition of
$W_V W_O$ is a rank-\texttt{d\_head} matrix in $d_{\text{model}} \times d_{\text{model}}$
space. The concepts in this section --- rank, transpose, identity, inverse --- are what let
you read those claims in papers without getting lost.
\end{motivbox}

\begin{exercise}[Rank by row reduction]
Find the rank of each matrix by reducing it to row echelon form.
\begin{enumerate}[label=(\alph*)]
    \item $A = \begin{bmatrix} 1 & 2 & 3 \\ 2 & 4 & 6 \\ 1 & 0 & 1 \end{bmatrix}$
    \item $B = \begin{bmatrix} 1 & 0 & 1 \\ 0 & 1 & 1 \\ 1 & 1 & 2 \end{bmatrix}$
    \item $C = \begin{bmatrix} 1 & 2 \\ 3 & 4 \\ 5 & 6 \end{bmatrix}$
          (a $3 \times 2$ matrix --- what is the maximum possible rank?)
\end{enumerate}
\end{exercise}

\begin{exercise}[Rank and the weight matrices]
In GPT-2 Small, each attention head has $d_{\text{model}} = 768$ and $d_{\text{head}} = 64$.
The query weight matrix $W_Q$ has shape $768 \times 64$.
\begin{enumerate}[label=(\alph*)]
    \item What is the maximum possible rank of $W_Q$?
    \item The QK matrix is defined as $W_{QK} = W_Q W_K^\top$. What shape is $W_{QK}$?
    \item What is the maximum possible rank of $W_{QK}$?
    \item Explain in plain language what it means for $W_{QK}$ to have rank 64 in a
          768-dimensional space. What are the heads ``blind'' to?
\end{enumerate}
\end{exercise}

\begin{exercise}[Transpose properties]
Let:
\[
    A = \begin{bmatrix} 1 & 2 & 3 \\ 4 & 5 & 6 \end{bmatrix}
\]
\begin{enumerate}[label=(\alph*)]
    \item Compute $A^\top$.
    \item Compute $A A^\top$ and $A^\top A$. What are their shapes?
    \item Is $A A^\top$ symmetric? (A matrix $M$ is symmetric if $M = M^\top$.)
          Is $A^\top A$ symmetric? State a general rule.
    \item In attention, the raw scores are $QK^\top$. If $Q$ has shape $(n \times d_h)$
          and $K$ has shape $(n \times d_h)$, what shape is $QK^\top$? Why is this
          the right shape for an attention pattern?
\end{enumerate}
\end{exercise}

\begin{exercise}[The identity and inverse]
\begin{enumerate}[label=(\alph*)]
    \item Write down the $3 \times 3$ identity matrix $I_3$.
    \item Verify that $I_3 \vec{v} = \vec{v}$ for $\vec{v} = \begin{bmatrix} 2 \\ -1 \\ 3 \end{bmatrix}$.
    \item Find the inverse of $A = \begin{bmatrix} 2 & 1 \\ 5 & 3 \end{bmatrix}$
          using the formula for $2 \times 2$ matrices:
          $A^{-1} = \frac{1}{ad - bc}\begin{bmatrix} d & -b \\ -c & a \end{bmatrix}$.
    \item Verify your answer: compute $AA^{-1}$ and confirm you get $I$.
    \item What is the determinant of $\begin{bmatrix} 2 & 4 \\ 1 & 2 \end{bmatrix}$?
          Does this matrix have an inverse? Why or why not?
\end{enumerate}
\end{exercise}

\begin{exercise}[Symmetric and orthogonal matrices]
\begin{enumerate}[label=(\alph*)]
    \item Which of the following are symmetric?
    \[
        P = \begin{bmatrix} 1 & 3 \\ 3 & 2 \end{bmatrix}, \quad
        Q = \begin{bmatrix} 0 & 1 \\ -1 & 0 \end{bmatrix}, \quad
        R = \begin{bmatrix} 4 & 1 & 2 \\ 1 & 0 & 3 \\ 2 & 3 & 5 \end{bmatrix}
    \]
    \item A matrix $U$ is \emph{orthogonal} if $U^\top U = I$.
          Verify that $U = \frac{1}{\sqrt{2}}\begin{bmatrix} 1 & -1 \\ 1 & 1 \end{bmatrix}$
          is orthogonal.
    \item If $U$ is orthogonal, what is $U^{-1}$? (You do not need to compute --- just
          read it from the definition.)
    \item Orthogonal matrices preserve vector lengths: $\|U\vec{v}\| = \|\vec{v}\|$.
          Verify this for $\vec{v} = \begin{bmatrix} 1 \\ 0 \end{bmatrix}$ and the
          matrix $U$ above.
\end{enumerate}
\end{exercise}

\begin{exercise}[Matrix multiplication order matters]
\begin{enumerate}[label=(\alph*)]
    \item Let $A = \begin{bmatrix} 1 & 2 \\ 0 & 1 \end{bmatrix}$ and
          $B = \begin{bmatrix} 1 & 0 \\ 3 & 1 \end{bmatrix}$.
          Compute $AB$ and $BA$. Are they equal?
    \item This is the non-commutativity of matrix multiplication: in general $AB \neq BA$.
          Give one concrete example of why the order of $W_Q, W_K$ matters in transformers.
    \item Compute $(AB)^\top$ and $B^\top A^\top$. Are they equal?
          This is the \emph{transpose of a product} rule: $(AB)^\top = B^\top A^\top$.
\end{enumerate}
\end{exercise}

\begin{exercise}[Trace and the identity]
The \emph{trace} of a square matrix $A$ is the sum of its diagonal entries:
$\mathrm{tr}(A) = \sum_i A_{ii}$.
\begin{enumerate}[label=(\alph*)]
    \item Compute $\mathrm{tr}(I_3)$.
    \item Compute $\mathrm{tr}(A)$ for $A = \begin{bmatrix} 3 & 1 & 0 \\ 2 & -1 & 4 \\ 0 & 2 & 5 \end{bmatrix}$.
    \item It can be shown that $\mathrm{tr}(AB) = \mathrm{tr}(BA)$ even though $AB \neq BA$ in general.
          Verify this for $A = \begin{bmatrix} 1 & 2 \\ 3 & 4 \end{bmatrix}$ and
          $B = \begin{bmatrix} 0 & 1 \\ 1 & 0 \end{bmatrix}$.
\end{enumerate}
\end{exercise}

% ──────────────────────────────────────────────────────────────────────────────
\section{Linear Mappings}
% ──────────────────────────────────────────────────────────────────────────────

\begin{motivbox}
\textbf{Why this matters.}
MML \S2.7 frames matrices not as arrays of numbers but as \emph{functions between vector
spaces}. This is the right way to think about transformer weight matrices: $W_Q$ is a
function that maps residual stream vectors to query space; $W_E$ (the embedding matrix) is a
function that maps token indices to vectors. The kernel and image of these functions tell you
what information they preserve, what they discard, and what they can express.
\end{motivbox}

\begin{exercise}[Verifying linearity]
A function $f: \mathbb{R}^n \to \mathbb{R}^m$ is \emph{linear} if:
(i) $f(\vec{u} + \vec{v}) = f(\vec{u}) + f(\vec{v})$, and
(ii) $f(\alpha \vec{v}) = \alpha f(\vec{v})$ for any scalar $\alpha$.

For each function below, determine whether it is linear. If not, give a counterexample.
\begin{enumerate}[label=(\alph*)]
    \item $f(x_1, x_2) = 3x_1 - 2x_2$
    \item $f(x_1, x_2) = x_1 x_2$
    \item $f(x_1, x_2) = x_1 + 1$
    \item $f(\vec{x}) = A\vec{x}$ for any fixed matrix $A$
\end{enumerate}
\end{exercise}

\begin{exercise}[Kernel (null space)]
The \emph{kernel} of a linear map $f$ is the set of all inputs that map to $\vec{0}$:
$\ker(f) = \{\vec{x} : f(\vec{x}) = \vec{0}\}$.

For $A = \begin{bmatrix} 1 & 2 & 3 \\ 2 & 4 & 6 \end{bmatrix}$:
\begin{enumerate}[label=(\alph*)]
    \item Solve $A\vec{x} = \vec{0}$ to find $\ker(A)$.
    \item What is the dimension of $\ker(A)$? (This is the \emph{nullity} of $A$.)
    \item The rank-nullity theorem states: $\mathrm{rank}(A) + \mathrm{nullity}(A) = n$,
          where $n$ is the number of columns. Verify this holds for $A$.
    \item In terms of transformer circuits: if $\ker(W_Q)$ is large, what does that mean
          about how much information the query projection ``throws away''?
\end{enumerate}
\end{exercise}

\begin{exercise}[Image (column space)]
The \emph{image} of a linear map $A$ is the set of all vectors $A$ can produce:
$\mathrm{im}(A) = \{A\vec{x} : \vec{x} \in \mathbb{R}^n\}$.

For $A = \begin{bmatrix} 1 & 0 \\ 0 & 1 \\ 1 & 1 \end{bmatrix}$:
\begin{enumerate}[label=(\alph*)]
    \item What is the shape of $A$? What are its domain and codomain?
    \item What is $\mathrm{rank}(A)$?
    \item The image of $A$ is spanned by the columns of $A$. Write out the two
          column vectors.
    \item Is the vector $\begin{bmatrix} 2 \\ 3 \\ 5 \end{bmatrix}$ in $\mathrm{im}(A)$?
          Is $\begin{bmatrix} 1 \\ 0 \\ 2 \end{bmatrix}$?
          (Check whether each can be written as $A\vec{x}$ for some $\vec{x}$.)
\end{enumerate}
\end{exercise}

\begin{exercise}[Composition of linear maps]
Let $f: \mathbb{R}^3 \to \mathbb{R}^2$ be given by matrix $A$, and
$g: \mathbb{R}^2 \to \mathbb{R}^2$ be given by matrix $B$:
\[
    A = \begin{bmatrix} 1 & 0 & 1 \\ 0 & 1 & 1 \end{bmatrix}, \qquad
    B = \begin{bmatrix} 2 & 0 \\ 0 & 3 \end{bmatrix}
\]
\begin{enumerate}[label=(\alph*)]
    \item Compute $C = BA$ (the matrix for $g \circ f$). What shape is $C$?
    \item Apply $f$ then $g$ to $\vec{v} = \begin{bmatrix} 1 \\ 2 \\ 1 \end{bmatrix}$
          step by step.
    \item Apply $C$ directly to $\vec{v}$ and confirm you get the same answer.
    \item In a transformer, the OV circuit $W_V W_O$ is a composition of two linear maps.
          $W_V$ maps from $d_{\text{model}}$ to $d_{\text{head}}$, and $W_O$ maps from
          $d_{\text{head}}$ back to $d_{\text{model}}$.
          What is the shape of $W_V W_O$? What is its maximum rank?
\end{enumerate}
\end{exercise}

\begin{exercise}[Isomorphism and dimension]
Two vector spaces are \emph{isomorphic} if there is a bijective linear map between them.
\begin{enumerate}[label=(\alph*)]
    \item Are $\mathbb{R}^2$ and $\mathbb{R}^3$ isomorphic? Why or why not?
          (Hint: think about what a bijective linear map would have to do to dimensions.)
    \item The space of $2 \times 2$ matrices $\mathbb{R}^{2 \times 2}$ has dimension 4.
          Is it isomorphic to $\mathbb{R}^4$? Describe an explicit bijection
          (a way to turn any $2 \times 2$ matrix into a 4-dimensional vector).
    \item In a transformer, the residual stream lives in $\mathbb{R}^{d_{\text{model}}}$.
          The query space for one head lives in $\mathbb{R}^{d_{\text{head}}}$.
          Are these isomorphic? What does $W_Q$ do to bridge them,
          and why is information necessarily lost?
\end{enumerate}
\end{exercise}

\begin{exercise}[Change of basis (conceptual)]
\begin{enumerate}[label=(\alph*)]
    \item The standard basis for $\mathbb{R}^2$ is
          $\vec{e}_1 = \begin{bmatrix} 1 \\ 0 \end{bmatrix}$,
          $\vec{e}_2 = \begin{bmatrix} 0 \\ 1 \end{bmatrix}$.
          Write $\vec{v} = \begin{bmatrix} 3 \\ -2 \end{bmatrix}$ as a linear combination
          of $\vec{e}_1$ and $\vec{e}_2$.
    \item Now suppose you want to express $\vec{v}$ in a new basis:
          $\vec{b}_1 = \begin{bmatrix} 1 \\ 1 \end{bmatrix}$,
          $\vec{b}_2 = \begin{bmatrix} 1 \\ -1 \end{bmatrix}$.
          Find scalars $c_1, c_2$ such that $c_1 \vec{b}_1 + c_2 \vec{b}_2 = \vec{v}$.
    \item In mechanistic interpretability, researchers often project residual stream
          activations onto specific directions (e.g., the ``unembedding direction'' of a
          token, or the top singular vectors of a weight matrix) to understand what
          information is present.
          How does this relate to the idea of expressing a vector in a different basis?
\end{enumerate}
\end{exercise}

% ──────────────────────────────────────────────────────────────────────────────
\section{Norms, Inner Products, and Orthogonality}
% ──────────────────────────────────────────────────────────────────────────────

\begin{motivbox}
\textbf{Why this matters.}
MML Chapter 3 (\S3.1--3.4) generalises the dot product to the concept of an \emph{inner product},
and introduces norms and orthogonality rigorously. In transformers: the $\sqrt{d_{\text{head}}}$
scaling in attention is a direct consequence of norm growth with dimension; orthogonality explains
why features in different subspaces don't interfere; and projection onto a direction is an inner
product operation. These exercises build the geometric intuition behind those mechanisms.
\end{motivbox}

\begin{exercise}[Computing norms]
The $\ell^2$ (Euclidean) norm of a vector $\vec{v} \in \mathbb{R}^n$ is
$\|\vec{v}\|_2 = \sqrt{\sum_i v_i^2}$.
\begin{enumerate}[label=(\alph*)]
    \item Compute $\|\vec{v}\|_2$ for $\vec{v} = \begin{bmatrix} 3 \\ -4 \end{bmatrix}$.
    \item Compute $\|\vec{u}\|_2$ for $\vec{u} = \begin{bmatrix} 1 \\ 1 \\ 1 \\ 1 \end{bmatrix}$.
    \item If each component of $\vec{u}$ is drawn from $\mathcal{N}(0,1)$ with $d$ components,
          the expected $\ell^2$ norm grows as $\sqrt{d}$. This is why attention logits are
          scaled by $\frac{1}{\sqrt{d_{\text{head}}}}$.
          For $d_{\text{head}} = 64$, what is the scaling factor?
    \item Compute the $\ell^1$ norm $\|\vec{v}\|_1 = \sum_i |v_i|$ for $\vec{v} = \begin{bmatrix} 3 \\ -4 \end{bmatrix}$.
          Is it larger or smaller than $\|\vec{v}\|_2$?
\end{enumerate}
\end{exercise}

\begin{exercise}[Inner products]
The standard dot product is one example of an inner product. A general inner product
$\langle \cdot, \cdot \rangle$ must be: (i) symmetric, (ii) linear in the first argument,
(iii) positive definite ($\langle \vec{v}, \vec{v} \rangle \geq 0$, with equality only if
$\vec{v} = \vec{0}$).
\begin{enumerate}[label=(\alph*)]
    \item Verify that the standard dot product $\langle \vec{u}, \vec{v} \rangle = \vec{u}^\top \vec{v}$
          satisfies positive definiteness for $\vec{v} = \begin{bmatrix} 2 \\ -3 \end{bmatrix}$.
    \item A \emph{weighted} inner product uses a positive definite matrix $A$:
          $\langle \vec{u}, \vec{v} \rangle_A = \vec{u}^\top A \vec{v}$.
          Compute $\langle \vec{u}, \vec{v} \rangle_A$ for
          $\vec{u} = \begin{bmatrix} 1 \\ 0 \end{bmatrix}$,
          $\vec{v} = \begin{bmatrix} 1 \\ 1 \end{bmatrix}$,
          $A = \begin{bmatrix} 2 & 0 \\ 0 & 3 \end{bmatrix}$.
    \item Is $f(\vec{u}, \vec{v}) = u_1 v_1 - u_2 v_2$ a valid inner product on $\mathbb{R}^2$?
          Check positive definiteness by trying $\vec{v} = \begin{bmatrix} 0 \\ 1 \end{bmatrix}$.
\end{enumerate}
\end{exercise}

\begin{exercise}[Orthogonality and projection]
Two vectors are \emph{orthogonal} if their inner product is zero: $\langle \vec{u}, \vec{v} \rangle = 0$.
\begin{enumerate}[label=(\alph*)]
    \item Are $\vec{u} = \begin{bmatrix} 1 \\ 2 \\ -1 \end{bmatrix}$ and
          $\vec{v} = \begin{bmatrix} 2 \\ 0 \\ 2 \end{bmatrix}$ orthogonal?
    \item The projection of $\vec{b}$ onto $\vec{a}$ is:
          $\mathrm{proj}_{\vec{a}} \vec{b} = \frac{\vec{a} \cdot \vec{b}}{\vec{a} \cdot \vec{a}} \vec{a}$
    Project $\vec{b} = \begin{bmatrix} 3 \\ 4 \end{bmatrix}$ onto $\vec{a} = \begin{bmatrix} 1 \\ 0 \end{bmatrix}$.
          Interpret the result geometrically.
    \item Compute the projection of $\vec{b} = \begin{bmatrix} 3 \\ 4 \end{bmatrix}$ onto
          $\vec{a} = \begin{bmatrix} 1 \\ 1 \end{bmatrix}$ (normalise $\vec{a}$ first).
    \item In mech interp, researchers project residual stream activations onto a
          ``token direction'' (e.g., the row of the unembedding matrix for token ``John'').
          This projection is an inner product. A high value means the residual stream is
          ``pointing toward'' that token's output direction. Does this feel concrete now?
\end{enumerate}
\end{exercise}

\begin{exercise}[Orthogonal complements]
The \emph{orthogonal complement} of a subspace $U$ in $\mathbb{R}^n$ is
$U^\perp = \{\vec{w} : \langle \vec{w}, \vec{u} \rangle = 0 \text{ for all } \vec{u} \in U\}$.
\begin{enumerate}[label=(\alph*)]
    \item Let $U = \mathrm{span}\!\left\{\begin{bmatrix}1\\1\\0\end{bmatrix}\right\}$ in $\mathbb{R}^3$.
          Find two vectors in $U^\perp$ (i.e., two independent vectors orthogonal to $\begin{bmatrix}1\\1\\0\end{bmatrix}$).
    \item What is $\dim(U) + \dim(U^\perp)$? What is the ambient dimension $n$? State the general rule.
    \item Every vector $\vec{x} \in \mathbb{R}^n$ can be uniquely decomposed as
          $\vec{x} = \vec{u} + \vec{u}^\perp$ where $\vec{u} \in U$ and $\vec{u}^\perp \in U^\perp$.
          Decompose $\vec{x} = \begin{bmatrix}2\\3\\1\end{bmatrix}$ into its component
          along $\vec{a} = \begin{bmatrix}1\\1\\0\end{bmatrix}$ and its orthogonal remainder.
    \item Why does orthogonal decomposition matter for low-rank matrices? If $W_Q$ has rank 64,
          what happens to residual stream components in $({\rm im}\, W_Q^\top)^\perp$?
\end{enumerate}
\end{exercise}

\begin{exercise}[Angles and the $\sqrt{d}$ scaling]
The angle $\theta$ between two vectors satisfies $\cos\theta = \frac{\vec{u} \cdot \vec{v}}{\|\vec{u}\|\|\vec{v}\|}$.
\begin{enumerate}[label=(\alph*)]
    \item Find the angle between $\vec{u} = \begin{bmatrix}1\\0\end{bmatrix}$ and
          $\vec{v} = \begin{bmatrix}1\\1\end{bmatrix}$ (your answer will be in radians or degrees).
    \item Suppose query vector $\vec{q}$ and key vector $\vec{k}$ both have unit norm
          and are at $45^\circ$ to each other. What is $\vec{q} \cdot \vec{k}$?
    \item Now suppose $d_{\text{head}} = 64$ and both $\vec{q}, \vec{k}$ have components
          drawn from $\mathcal{N}(0, 1)$, so their expected norm is $\sqrt{64} = 8$.
          What is the expected scale of $\vec{q} \cdot \vec{k}$ if the angle is still $45^\circ$?
    \item After dividing by $\sqrt{d_{\text{head}}} = 8$, what is the scaled attention logit?
          Explain in one sentence why this scaling is necessary for softmax to produce
          well-spread attention weights rather than near-one-hot distributions.
\end{enumerate}
\end{exercise}

% ──────────────────────────────────────────────────────────────────────────────
\newpage
\section*{Solutions}
\label{sec:solutions}
\addcontentsline{toc}{section}{Solutions}
% ──────────────────────────────────────────────────────────────────────────────

\setcounter{section}{1}
\setcounter{solution}{0}

% ── Section 1 Solutions ───────────────────────────────────────────────────────
\subsection*{Section 1: Systems of Linear Equations}

\begin{solution}[$2 \times 2$ system]
\begin{enumerate}[label=(\alph*)]
    \item Augmented matrix: $\begin{bmatrix} 2 & 1 & \mid & 5 \\ 1 & -1 & \mid & 1 \end{bmatrix}$
    \item $R_1 \leftrightarrow R_2$: $\begin{bmatrix} 1 & -1 & \mid & 1 \\ 2 & 1 & \mid & 5 \end{bmatrix}$,
          then $R_2 \leftarrow R_2 - 2R_1$:
          $\begin{bmatrix} 1 & -1 & \mid & 1 \\ 0 & 3 & \mid & 3 \end{bmatrix}$
    \item $y = 1$, back-substitute: $x - 1 = 1$, so $x = 2$. Solution: $(x, y) = (2, 1)$.
\end{enumerate}
\end{solution}

\begin{solution}[Gaussian elimination on a $3 \times 3$ system]
\begin{enumerate}[label=(\alph*)]
    \item $\begin{bmatrix} 1 & 2 & 1 & \mid & 6 \\ 2 & 1 & 3 & \mid & 9 \\ 0 & 1 & 1 & \mid & 3 \end{bmatrix}$
    \item $R_2 \leftarrow R_2 - 2R_1$:
          $\begin{bmatrix} 1 & 2 & 1 & \mid & 6 \\ 0 & -3 & 1 & \mid & -3 \\ 0 & 1 & 1 & \mid & 3 \end{bmatrix}$
    \item $R_3 \leftarrow R_3 + \frac{1}{3}R_2$:
          $\begin{bmatrix} 1 & 2 & 1 & \mid & 6 \\ 0 & -3 & 1 & \mid & -3 \\ 0 & 0 & \frac{4}{3} & \mid & 2 \end{bmatrix}$
    \item $x_3 = 2 \cdot \frac{3}{4} = \frac{3}{2}$. From row 2: $-3x_2 + \frac{3}{2} = -3$, so $x_2 = \frac{3}{2}$.
          From row 1: $x_1 + 3 + \frac{3}{2} = 6$, so $x_1 = \frac{3}{2}$.
\end{enumerate}
\end{solution}

\begin{solution}[Counting solutions]
\begin{enumerate}[label=(\alph*)]
    \item \textbf{One solution.} Two equations, two unknowns, no contradiction. Back-substitute: $x_2 = 1$, $x_1 = 1$.
    \item \textbf{No solution.} Row 2 says $0 = 5$, which is a contradiction.
    \item \textbf{Infinitely many solutions.} The last row is all zeros --- the system is underdetermined.
          One free variable ($x_3$); solutions form a 1-dimensional family.
\end{enumerate}
\end{solution}

\begin{solution}[Homogeneous systems]
\begin{enumerate}[label=(\alph*)]
    \item $\vec{x} = \vec{0}$ is always a solution (the \emph{trivial solution}).
    \item Row reduce:
          $\begin{bmatrix} 1 & 2 & -1 \\ 2 & 4 & -2 \end{bmatrix} \to
           \begin{bmatrix} 1 & 2 & -1 \\ 0 & 0 & 0 \end{bmatrix}$
          Row 2 is redundant. Free variables: $x_2 = s$, $x_3 = t$. Then $x_1 = -2s + t$.
          General solution: $\begin{bmatrix} -2s + t \\ s \\ t \end{bmatrix}
          = s\begin{bmatrix} -2 \\ 1 \\ 0 \end{bmatrix}
          + t\begin{bmatrix} 1 \\ 0 \\ 1 \end{bmatrix}$
    \item Dimension 2. The null space is spanned by two independent vectors --- the matrix ``throws away''
          a 2-dimensional subspace of its input.
\end{enumerate}
\end{solution}

\begin{solution}[Geometric interpretation of solutions]
\begin{enumerate}[label=(\alph*)]
    \item (i) Lines intersect → unique solution. (ii) Parallel, distinct → no solution. (iii) Same line → infinitely many solutions.
    \item $2x + 2y = 6$ is $2 \times (x + y = 3)$ --- the same line. Infinitely many solutions:
          $\{(t, 3-t) : t \in \mathbb{R}\}$. A whole line of solutions.
    \item Row reduce: $R_2 - 2R_1$ gives $0 = 1$ --- no solution. Parallel lines.
\end{enumerate}
\end{solution}

\begin{solution}[Back-substitution practice]
From row 3: $2x_3 = 6 \Rightarrow x_3 = 3$.
From row 2: $3x_2 - 3 = 4 \Rightarrow x_2 = \frac{7}{3}$.
From row 1: $2x_1 - \frac{7}{3} + 9 = 7 \Rightarrow 2x_1 = 7 + \frac{7}{3} - 9 = \frac{1}{3} \Rightarrow x_1 = \frac{1}{6}$.
\end{solution}

\setcounter{section}{2}
\setcounter{solution}{0}

% ── Section 2 Solutions ───────────────────────────────────────────────────────
\newpage
\subsection*{Section 2: Matrix Properties}

\begin{solution}[Rank by row reduction]
\begin{enumerate}[label=(\alph*)]
    \item $R_2 \leftarrow R_2 - 2R_1$, $R_3 \leftarrow R_3 - R_1$:
          $\begin{bmatrix} 1 & 2 & 3 \\ 0 & 0 & 0 \\ 0 & -2 & -2 \end{bmatrix}$,
          swap $R_2 \leftrightarrow R_3$: two nonzero rows. $\mathrm{rank}(A) = 2$.
    \item $R_3 \leftarrow R_3 - R_1 - R_2$:
          $\begin{bmatrix} 1 & 0 & 1 \\ 0 & 1 & 1 \\ 0 & 0 & 0 \end{bmatrix}$. $\mathrm{rank}(B) = 2$.
    \item Maximum possible rank is $\min(3, 2) = 2$. After reduction: two independent rows. $\mathrm{rank}(C) = 2$.
\end{enumerate}
\end{solution}

\begin{solution}[Rank and the weight matrices]
\begin{enumerate}[label=(\alph*)]
    \item $\min(768, 64) = 64$.
    \item $W_Q$ is $(768 \times 64)$, $W_K^\top$ is $(64 \times 768)$, so $W_{QK}$ is $(768 \times 768)$.
    \item At most $64$ (bounded by the intermediate dimension).
    \item The head lives in a 64-dimensional subspace of a 768-dimensional space. Any residual stream
          direction orthogonal to this subspace is invisible to the head's attention pattern --- it
          contributes nothing to the query or key computation.
\end{enumerate}
\end{solution}

\begin{solution}[Transpose properties]
\begin{enumerate}[label=(\alph*)]
    \item $A^\top = \begin{bmatrix} 1 & 4 \\ 2 & 5 \\ 3 & 6 \end{bmatrix}$
    \item $AA^\top$ is $(2 \times 3)(3 \times 2) = (2 \times 2)$;
          $A^\top A$ is $(3 \times 2)(2 \times 3) = (3 \times 3)$.
    \item Both are symmetric. General rule: $AA^\top$ and $A^\top A$ are always symmetric for any $A$.
    \item $Q$ is $(n \times d_h)$, $K^\top$ is $(d_h \times n)$, so $QK^\top$ is $(n \times n)$ --- one
          score per (query position, key position) pair, which is exactly an attention weight matrix.
\end{enumerate}
\end{solution}

\begin{solution}[The identity and inverse]
\begin{enumerate}[label=(\alph*)]
    \item $I_3 = \begin{bmatrix} 1 & 0 & 0 \\ 0 & 1 & 0 \\ 0 & 0 & 1 \end{bmatrix}$
    \item $I_3 \vec{v} = \begin{bmatrix} 2 \\ -1 \\ 3 \end{bmatrix}$ (unchanged).
    \item $\det(A) = 2 \cdot 3 - 1 \cdot 5 = 1$. So $A^{-1} = \begin{bmatrix} 3 & -1 \\ -5 & 2 \end{bmatrix}$.
    \item $AA^{-1} = \begin{bmatrix} 2 & 1 \\ 5 & 3 \end{bmatrix}\begin{bmatrix} 3 & -1 \\ -5 & 2 \end{bmatrix}
          = \begin{bmatrix} 1 & 0 \\ 0 & 1 \end{bmatrix}$. \checkmark
    \item $\det = 2 \cdot 2 - 4 \cdot 1 = 0$. No inverse --- the matrix is \emph{singular} (rank 1, not full rank).
\end{enumerate}
\end{solution}

\begin{solution}[Symmetric and orthogonal matrices]
\begin{enumerate}[label=(\alph*)]
    \item $P$ and $R$ are symmetric ($P^\top = P$, $R^\top = R$). $Q$ is not: $Q^\top = \begin{bmatrix} 0 & -1 \\ 1 & 0 \end{bmatrix} \neq Q$.
    \item $U^\top U = \frac{1}{2}\begin{bmatrix} 1 & 1 \\ -1 & 1 \end{bmatrix}
          \begin{bmatrix} 1 & -1 \\ 1 & 1 \end{bmatrix}
          = \frac{1}{2}\begin{bmatrix} 2 & 0 \\ 0 & 2 \end{bmatrix} = I$. \checkmark
    \item $U^{-1} = U^\top$ (read directly from $U^\top U = I$).
    \item $U\vec{v} = \frac{1}{\sqrt{2}}\begin{bmatrix} 1 \\ 1 \end{bmatrix}$.
          $\|U\vec{v}\| = \frac{1}{\sqrt{2}}\sqrt{1^2 + 1^2} = \frac{\sqrt{2}}{\sqrt{2}} = 1 = \|\vec{v}\|$. \checkmark
\end{enumerate}
\end{solution}

\begin{solution}[Matrix multiplication order matters]
\begin{enumerate}[label=(\alph*)]
    \item $AB = \begin{bmatrix}7&2\\3&1\end{bmatrix}$, $BA = \begin{bmatrix}1&2\\3&7\end{bmatrix}$. Not equal.
    \item The QK matrix $W_Q W_K^\top$ defines the attention pattern; $W_K^\top W_Q$ has a different shape
          and different meaning entirely. Order encodes the direction of the mapping.
    \item $(AB)^\top = \begin{bmatrix}7&3\\2&1\end{bmatrix}$.
          $B^\top A^\top = \begin{bmatrix}1&3\\0&1\end{bmatrix}\begin{bmatrix}1&0\\2&1\end{bmatrix} = \begin{bmatrix}7&3\\2&1\end{bmatrix}$. Equal. \checkmark
\end{enumerate}
\end{solution}

\begin{solution}[Trace and the identity]
\begin{enumerate}[label=(\alph*)]
    \item $\mathrm{tr}(I_3) = 1 + 1 + 1 = 3$.
    \item $\mathrm{tr}(A) = 3 + (-1) + 5 = 7$.
    \item $AB = \begin{bmatrix}2&1\\4&3\end{bmatrix}$, $\mathrm{tr}(AB) = 5$.
          $BA = \begin{bmatrix}3&4\\1&2\end{bmatrix}$, $\mathrm{tr}(BA) = 5$. \checkmark
\end{enumerate}
\end{solution}

\setcounter{section}{3}
\setcounter{solution}{0}

% ── Section 3 Solutions ───────────────────────────────────────────────────────
\newpage
\subsection*{Section 3: Linear Mappings}

\begin{solution}[Verifying linearity]
\begin{enumerate}[label=(\alph*)]
    \item \textbf{Linear.} $f(\vec{u} + \vec{v}) = 3(u_1 + v_1) - 2(u_2 + v_2) = f(\vec{u}) + f(\vec{v})$;
          $f(\alpha\vec{v}) = \alpha f(\vec{v})$. Both hold.
    \item \textbf{Not linear.} $f(2\vec{v}) = 2v_1 \cdot 2v_2 = 4v_1 v_2 \neq 2 f(\vec{v})$.
    \item \textbf{Not linear.} $f(\vec{0}) = 0 + 1 = 1 \neq 0$. Linear maps must send $\vec{0}$ to $\vec{0}$.
    \item \textbf{Linear.} Matrix multiplication satisfies both linearity conditions for any fixed $A$.
\end{enumerate}
\end{solution}

\begin{solution}[Kernel (null space)]
\begin{enumerate}[label=(\alph*)]
    \item Row reduce: $R_2 \leftarrow R_2 - 2R_1$:
          $\begin{bmatrix} 1 & 2 & 3 \\ 0 & 0 & 0 \end{bmatrix}$.
          Free variables: $x_2 = s$, $x_3 = t$. Then $x_1 = -2s - 3t$.
          $\ker(A) = \mathrm{span}\!\left\{\begin{bmatrix} -2 \\ 1 \\ 0 \end{bmatrix},
          \begin{bmatrix} -3 \\ 0 \\ 1 \end{bmatrix}\right\}$
    \item Nullity $= 2$.
    \item $\mathrm{rank}(A) = 1$ (one nonzero row). $1 + 2 = 3 = $ number of columns. \checkmark
    \item A large kernel means the query projection discards a large subspace of the residual stream ---
          many residual stream directions contribute nothing to the head's attention pattern.
\end{enumerate}
\end{solution}

\begin{solution}[Image (column space)]
\begin{enumerate}[label=(\alph*)]
    \item Shape: $3 \times 2$. Domain: $\mathbb{R}^2$, codomain: $\mathbb{R}^3$.
    \item Rank 2 (columns are independent).
    \item Column vectors: $\begin{bmatrix} 1 \\ 0 \\ 1 \end{bmatrix}$ and $\begin{bmatrix} 0 \\ 1 \\ 1 \end{bmatrix}$.
    \item For $\begin{bmatrix} 2 \\ 3 \\ 5 \end{bmatrix}$: need $x_1 = 2$, $x_2 = 3$, check third component $2 + 3 = 5$. \checkmark In the image.
          For $\begin{bmatrix} 1 \\ 0 \\ 2 \end{bmatrix}$: need $x_1 = 1$, $x_2 = 0$, check $1 + 0 = 1 \neq 2$. \textbf{Not} in the image.
\end{enumerate}
\end{solution}

\begin{solution}[Composition of linear maps]
\begin{enumerate}[label=(\alph*)]
    \item $C = BA = \begin{bmatrix} 2 & 0 \\ 0 & 3 \end{bmatrix}\begin{bmatrix} 1 & 0 & 1 \\ 0 & 1 & 1 \end{bmatrix}
          = \begin{bmatrix} 2 & 0 & 2 \\ 0 & 3 & 3 \end{bmatrix}$. Shape: $2 \times 3$.
    \item $f(\vec{v}) = \begin{bmatrix} 1+1 \\ 2+1 \end{bmatrix} = \begin{bmatrix} 2 \\ 3 \end{bmatrix}$.
          $g(f(\vec{v})) = \begin{bmatrix} 4 \\ 9 \end{bmatrix}$.
    \item $C\vec{v} = \begin{bmatrix} 2+0+2 \\ 0+6+3 \end{bmatrix} = \begin{bmatrix} 4 \\ 9 \end{bmatrix}$. \checkmark
    \item $W_V$ is $(d_h \times d_m)$, $W_O$ is $(d_m \times d_h)$... Actually per MML convention with row vectors:
          $W_V$ maps $(d_m \to d_h)$ and $W_O$ maps $(d_h \to d_m)$, so $W_V W_O$ is $(d_m \times d_m)$.
          Maximum rank: $d_h = 64$.
\end{enumerate}
\end{solution}

\begin{solution}[Isomorphism and dimension]
\begin{enumerate}[label=(\alph*)]
    \item No. A bijective linear map preserves dimension, but $\mathbb{R}^2$ has dimension 2 and
          $\mathbb{R}^3$ has dimension 3. No such bijection exists.
    \item Yes. An explicit bijection: $\begin{bmatrix}a&b\\c&d\end{bmatrix} \mapsto \begin{bmatrix}a\\b\\c\\d\end{bmatrix}$.
          This is linear, injective, and surjective.
    \item $\mathbb{R}^{768}$ and $\mathbb{R}^{64}$ are not isomorphic (different dimensions).
          $W_Q$ maps from the larger space to the smaller one --- it is a surjective (onto)
          but not injective map. Its kernel has dimension $768 - 64 = 704$: a 704-dimensional
          subspace of the residual stream maps to zero and is completely invisible to the head.
\end{enumerate}
\end{solution}

\begin{solution}[Change of basis]
\begin{enumerate}[label=(\alph*)]
    \item $\vec{v} = 3\vec{e}_1 + (-2)\vec{e}_2$.
    \item Solve $c_1 + c_2 = 3$, $c_1 - c_2 = -2$. Adding: $2c_1 = 1$, so $c_1 = \frac{1}{2}$, $c_2 = \frac{5}{2}$.
    \item Projecting a residual stream vector onto a specific direction (e.g., the ``John'' direction in the IOI
          circuit) is exactly expressing that vector's coordinates in a new basis where that direction is
          one of the basis vectors. The projection coefficient tells you ``how much'' of that concept is present.
\end{enumerate}
\end{solution}

\setcounter{section}{4}
\setcounter{solution}{0}

% ── Section 4 Solutions ───────────────────────────────────────────────────────
\newpage
\subsection*{Section 4: Norms, Inner Products, and Orthogonality}

\begin{solution}[Computing norms]
\begin{enumerate}[label=(\alph*)]
    \item $\|\vec{v}\|_2 = \sqrt{9 + 16} = \sqrt{25} = 5$.
    \item $\|\vec{u}\|_2 = \sqrt{1+1+1+1} = 2$.
    \item $\frac{1}{\sqrt{64}} = \frac{1}{8} = 0.125$.
    \item $\|\vec{v}\|_1 = 3 + 4 = 7$. Larger than $\|\vec{v}\|_2 = 5$. In general $\|\vec{v}\|_1 \geq \|\vec{v}\|_2$.
\end{enumerate}
\end{solution}

\begin{solution}[Inner products]
\begin{enumerate}[label=(\alph*)]
    \item $\langle \vec{v}, \vec{v} \rangle = 4 + 9 = 13 > 0$. \checkmark
    \item $\vec{u}^\top A \vec{v} = \begin{bmatrix}1&0\end{bmatrix}\begin{bmatrix}2&0\\0&3\end{bmatrix}\begin{bmatrix}1\\1\end{bmatrix}
          = \begin{bmatrix}2&0\end{bmatrix}\begin{bmatrix}1\\1\end{bmatrix} = 2$.
    \item $f(\vec{v},\vec{v}) = 0 \cdot 0 - 1 \cdot 1 = -1 < 0$. Not positive definite --- not a valid inner product.
\end{enumerate}
\end{solution}

\begin{solution}[Orthogonality and projection]
\begin{enumerate}[label=(\alph*)]
    \item $\vec{u} \cdot \vec{v} = 2 + 0 - 2 = 0$. Yes, orthogonal.
    \item $\mathrm{proj}_{\vec{a}}\vec{b} = \frac{3}{1}\begin{bmatrix}1\\0\end{bmatrix} = \begin{bmatrix}3\\0\end{bmatrix}$.
          This is the shadow of $\vec{b}$ onto the $x$-axis --- the $x$-component of $\vec{b}$.
    \item $\hat{\vec{a}} = \frac{1}{\sqrt{2}}\begin{bmatrix}1\\1\end{bmatrix}$.
          $\hat{\vec{a}} \cdot \vec{b} = \frac{1}{\sqrt{2}}(3+4) = \frac{7}{\sqrt{2}}$.
          Projection: $\frac{7}{\sqrt{2}} \cdot \frac{1}{\sqrt{2}}\begin{bmatrix}1\\1\end{bmatrix}
          = \frac{7}{2}\begin{bmatrix}1\\1\end{bmatrix} = \begin{bmatrix}3.5\\3.5\end{bmatrix}$.
    \item Yes --- the logit lens and direct path analysis are both inner product operations.
\end{enumerate}
\end{solution}

\begin{solution}[Orthogonal complements]
\begin{enumerate}[label=(\alph*)]
    \item Need $w_1 + w_2 = 0$. Two independent solutions:
          $\begin{bmatrix}1\\-1\\0\end{bmatrix}$ and $\begin{bmatrix}0\\0\\1\end{bmatrix}$.
    \item $\dim(U) = 1$, $\dim(U^\perp) = 2$, sum $= 3 = n$. General rule: $\dim(U) + \dim(U^\perp) = n$.
    \item Project onto $\hat{\vec{a}} = \frac{1}{\sqrt{2}}\begin{bmatrix}1\\1\\0\end{bmatrix}$:
          coefficient $= \frac{1}{\sqrt{2}}(2+3+0) = \frac{5}{\sqrt{2}}$.
          Component along $\vec{a}$: $\frac{5}{2}\begin{bmatrix}1\\1\\0\end{bmatrix} = \begin{bmatrix}2.5\\2.5\\0\end{bmatrix}$.
          Remainder: $\begin{bmatrix}2\\3\\1\end{bmatrix} - \begin{bmatrix}2.5\\2.5\\0\end{bmatrix} = \begin{bmatrix}-0.5\\0.5\\1\end{bmatrix}$.
    \item Those components are completely invisible to the query projection --- they contribute nothing
          to attention scores and are effectively discarded by the head.
\end{enumerate}
\end{solution}

\begin{solution}[Angles and the $\sqrt{d}$ scaling]
\begin{enumerate}[label=(\alph*)]
    \item $\cos\theta = \frac{0+0}{\sqrt{1}\sqrt{2}} \cdot \ldots$
          Actually: $\vec{u}\cdot\vec{v} = 1$, $\|\vec{u}\| = 1$, $\|\vec{v}\| = \sqrt{2}$.
          $\cos\theta = \frac{1}{\sqrt{2}}$, so $\theta = 45^\circ$.
    \item $\vec{q}\cdot\vec{k} = \cos(45^\circ) = \frac{1}{\sqrt{2}} \approx 0.707$.
    \item Each has expected norm $8$; $\vec{q}\cdot\vec{k} \approx \|\vec{q}\|\|\vec{k}\|\cos\theta
          = 8 \cdot 8 \cdot \frac{1}{\sqrt{2}} \approx 45$.
    \item After scaling: $45 / 8 \approx 5.7$. Without scaling, logits of magnitude $\sim 45$
          push softmax into saturation (near one-hot), making gradients vanish and the model
          unable to attend to multiple positions simultaneously.
\end{enumerate}
\end{solution}

\end{document}
